% Part: second-order-logic
% Chapter: syntax-and-semantics
% Section: semantic-notions

\documentclass[../../../include/open-logic-section]{subfiles}

\begin{document}

\olfileid[tr]{sol}{syn}{sem}
\olsection{Semantik Kavramlar}

\begin{explain}
\emph{Geçerlilik}, \emph{gerektirme} ve \emph{sağlanabilirlik} temel
mantıksal kavramları, ikinci dereceden mantık için birinci dereceden
mantıktakiyle aynı biçimde tanımlanır; tek fark, temeldeki sağlama
bağıntısının artık ikinci dereceden !!{formula}s için olan bağıntı olmasıdır.
Elbette ikinci dereceden bir !!{sentence}, yüklem ve fonksiyon değişkenleri
de içinde olmak üzere bütün değişkenlerin bağlı olduğu !!a{formula}{}dür.
\end{explain}

\begin{defn}[Geçerlilik]
Bir $!A$ tümcesi, ancak ve ancak her !!{structure}~$\Struct M$ için
$\Sat{M}{!A}$ ise \emph{geçerlidir}; bunu $\Entails !A$ ile gösteririz.
\end{defn}

\begin{defn}[Gerektirme]
Bir tümceler kümesi~$\Gamma$, ancak ve ancak $\Sat{M}{\Gamma}$ olan her
!!{structure}~$\Struct M$ için $\Sat{M}{!A}$ ise $!A$ tümcesini
\emph{gerektirir}; bunu $\Gamma \Entails !A$ ile gösteririz.
\end{defn}

\begin{defn}[Sağlanabilirlik]
Bir tümceler kümesi~$\Gamma$, bazı !!{structure}~$\Struct M$ için
$\Sat{M}{\Gamma}$ ise \emph{sağlanabilir}dir. $\Gamma$ sağlanabilir değilse
ona \emph{sağlanamaz} denir.
\end{defn}

\end{document}
